
% ---------- IM364TA Human Resource Management & Analytics (Improvement CIE)
\begin{iempaper}
\paperinfo{shownote=0,date={19\textsuperscript{th} June 2026},exam={Improvement CIE},maxmarks={10 + 50},sem={VI},prog={UG},duration={30 + 90 Min},
 title={Human Resource Management \& Analytics},code={IM364TA}}
\iemhead
\ptitle{Part -- A}
\begin{longtable}{|C{1.05cm}|p{12.2cm}|C{0.95cm}|C{0.95cm}|C{0.95cm}|}
\hline \textbf{Sl.\newline No.} & \multicolumn{1}{c|}{\textbf{Questions}} & \textbf{M} & \textbf{BT} & \textbf{CO}\\\hline
\q{1.}{Differentiate between HR Metrics and HR Analytics.}{02}{L2}{CO4}
\q{2.}{State any two objectives of an HR Scorecard.}{02}{L1}{CO4}
\q{3.}{Identify any two-recruitment metrics used in HR Analytics.}{02}{L1}{CO4}
\q{4.}{List any two benefits of HR Dashboards.}{02}{L2}{CO4}
\q{5.}{Name any two phases of HR Predictive Modelling.}{02}{L1}{CO4}
\end{longtable}
\ptitle{Part -- B}
\begin{longtable}{|C{1.05cm}|p{12.2cm}|C{0.95cm}|C{0.95cm}|C{0.95cm}|}
\hline \textbf{Sl.\newline No.} & \multicolumn{1}{c|}{\textbf{Questions}} & \textbf{M} & \textbf{BT} & \textbf{CO}\\\hline
\q{1.}{Explain the alignment of HR Analytics with organizational business goals and strategies. Illustrate with a suitable example.}{10}{L3}{CO4}
\q{2.}{Discuss the HR Analytics Framework and explain how it supports effective HR decision-making.}{10}{L3}{CO4}
\q{3.}{Discuss the role of HRIS in supporting HR decision-making and its applications in recruitment, training, and performance management.}{10}{L3}{CO4}
\q{4.}{Design an HR Dashboard for monitoring recruitment effectiveness. Identify suitable metrics and justify their inclusion.}{10}{L4}{CO4}
\q{5.}{A technology company has experienced an employee attrition rate of 18\% during the past year, significantly higher than the industry average. Available HR data indicate declining employee engagement scores, increased overtime hours, and reduced participation in training programs. Using the phases of HR Predictive Modelling, outline an approach to predict future attrition and recommend suitable retention strategies.}{10}{L4}{CO4}
\end{longtable}
\end{iempaper}

% ---------- IM365TDA Facilities Planning and Design (Improvement CIE)
\begin{iempaper}
\paperinfo{shownote=0,date={18\textsuperscript{th} June 2026},exam={Improvement CIE},maxmarks={10 + 50},sem={VI},prog={UG},duration={30 + 90 Min},
 title={Facilities Planning And Design},code={IM365TDA}}
\iemhead
{\small Note:}\par\vspace{-1pt}
\hspace*{0.5cm}{\footnotesize 1. \hspace{0.1cm}Answer all the Questions.}\par\vspace{3pt}
\begin{cietable}{M}{BT}{CO}
\qpart{Part -- A}
\q{1.}{List any two characteristics of fixed automation systems.}{(02)}{1}{1}
\q{2.}{What is Flexible Manufacturing System (FMS)?}{(02)}{2}{2}
\q{3.}{State the concept of Just-In-Time (JIT) manufacturing.}{(02)}{3}{3}
\q{4.}{Why is a Flexible Manufacturing System (FMS) more suitable than fixed automation for product variety?}{(02)}{2}{4}
\q{5.}{Why is fixed automation not suitable for low-volume customized products? Give two reasons}{(02)}{2}{3}
\qpart{Part -- B}
\q{1.}{Discuss characteristics and applications in manufacturing industries of fixed automation systems.}{(10)}{2}{2}
\q{2.}{Describe the working of a Flexible Manufacturing System (FMS).}{(10)}{3}{2}
\q{3.}{Explain the operation of a single-stage multi-machine system with a neat diagram.}{(10)}{3}{3}
\q{4.}{Illustrate how Just-In-Time (JIT) manufacturing contributes to waste elimination and productivity improvement. Explain its key principles.}{(10)}{3}{4}
\q{5.}{Propose a comprehensive plan to transform a traditional manufacturing facility into a lean production system based on JIT concepts.}{(10)}{2}{4}
\end{cietable}
\stars{********}
\end{iempaper}

% ---------- IM365TDB Service Operations Management (CIE - III)
\begin{iempaper}
\paperinfo{shownote=0,date={17-06-2026},exam={CIE -- III},maxmarks={10 + 50},sem={VI},prog={UG},duration={30 + 90 Min},
 title={Service Operations Management},code={IM365TDB}}
\iemhead
{\small Note:}\par\vspace{-1pt}
\hspace*{0.5cm}{\footnotesize Answer all the Questions.}\par\vspace{3pt}
\begin{cietable}{M}{BT}{CO}
\qpart{Part -- A}
\q{1.}{What is `mass customisation' and which two principles underpin it in services?}{2}{1}{2}
\q{2.}{What is the `market-operations gap' and what risk does it pose for service organisations?}{2}{2}{3}
\q{3.}{When a bank introduces a mobile app for all routine transfers, moving the labor from the teller to the user, what role is the customer now playing?}{2}{2}{2}
\q{4.}{What is meant by customer participation in services?}{2}{2}{4}
\q{5.}{A fast-food chain uses average service time per customer as its only performance measure. What is the risk?}{2}{2}{3}
\qpart{Part -- B}
\q{1}{Discuss the Hackman and Oldham Job Characteristics Model in the context of motivating service providers. How do specific job dimensions lead to better psychological states?}{10}{2}{4}
\q{2.}{Explain the different types of service processes and their importance in service operations management.}{10}{2}{4}
\q{3.}{Prestige Home Services (PHS) is a premium home maintenance company offering plumbing, electrical, carpentry, and deep-cleaning services in Hyderabad. Technicians are dispatched to customer homes and are expected to assess the problem, carry out repairs, and collect payment - all independently. PHS recently won a large corporate contract requiring 200 service visits per month across office premises. To meet this demand, PHS hired 40 new technicians rapidly with minimal on boarding. Within 3 months, customer complaints surged. Investigations revealed that technicians were arriving without knowing the scope of the job, improvising repairs without proper tools, quoting inconsistent prices to customers, and in several cases behaving rudely when customers questioned their work. Experienced technicians reported feeling demoralised as their reputation was being damaged by new hires. The operations manager is tasked with redesigning both the service process and the people management system before the corporate contract renewal in 4 months. (a) Using service blueprinting, redesign the end-to-end service process for PHS, identifying customer actions, front-stage actions, back-stage actions, support processes, and key fail points with solutions.}{10}{3}{3}
\end{cietable}
\end{iempaper}

% ---------- IM365TDD Design of Experiments (CIE Improvement)
\begin{iempaper}
\paperinfo{shownote=0,date={18\textsuperscript{th} JUNE 2026},exam={CIE Improvement},maxmarks={10 + 50},sem={VI},prog={UG},duration={30 + 90 Min},
 title={Design of Experiments},code={IM365TDD}}
\iemhead
{\small Instructions to Students:}\par\vspace{-1pt}
\hspace*{0.5cm}{\footnotesize 1. \hspace{0.1cm}Answer all questions}\par
\hspace*{0.5cm}{\footnotesize 2. \hspace{0.1cm}Use of Statistical tables and Orthogonal Arrays tables permitted}\par\vspace{3pt}
\begin{cietable}{M}{BT}{CO}
\qpart{Part -- A}
\q{1.}{What is the primary purpose of a ``Linear Graph'' in Taguchi's orthogonal array designs?}{2}{2}{CO3}
\q{2.}{What is the ``Strategy for Constructing an Orthogonal Array'' when the number of factors is high but the budget is low.}{2}{3}{CO3}
\q{3.}{Illustrate an application of the Dummy Level Technique for a factor with 2 levels in an L$_9$ (3-level) array.}{2}{2}{CO3}
\q{4.}{What is the primary purpose of a ``Verification Experiment'' in the robust design cycle?}{2}{2}{CO2}
\q{5.}{Contrast ``Testing Conditions'' with ``Control Factor Levels'' in a typical study.}{2}{2}{CO2}
\qpart{Part -- B}
\q{1}{In an investigation on the stability of power factor of capacitor it was decided to examine the following factors:
\qtab{|c|l|c|c|}{\hline \textbf{Sl No} & \multicolumn{1}{c|}{\textbf{Factor}} & \multicolumn{2}{c|}{\textbf{Levels}}\\\hline 1 & Leakage Current of Foil (A) & $\leq$ 2 mA & $\geq$ 2 mA\\\hline 2 & Power Factor of Foil (B) & $\leq$ 3\% & $\geq$ 3\%\\\hline 3 & Capacitance Value of Coil \textcopyright & 8-10 $\mu$f & 11-13 $\mu$f\\\hline 4 & Soaking Time during Impregnation & 2 hrs & 4 hrs\\\hline 5 & Loading Pattern & 70 \% & 100 \%\\\hline}
Suspected interaction effects are AB, BC. Develop the Orthogonal Array Layout and experimental plan.}{15}{4}{CO3}
\q{2}{In a study to examine process control parameter for joining ``BUS-BAR'' of an electrical system the following factors are levels were identified for investigation with respect to efficiency of joining
\qtabc{|c|l|c|c|}{\hline \multirow{2}{*}{\textbf{Factor Code}} & \multirow{2}{*}{\textbf{Factor Description}} & \multicolumn{2}{c|}{\textbf{Level}}\\\cline{3-4} & & 1 & 2\\\hline A & Torque & 25 ft lb & 40 ft lb\\\hline B & Compound & Nil & Silicon Grease\\\hline C & Surface Finish & As is & Ground\\\hline D & Joint Overlap & 8 mm & 16 mm\\\hline}
Develop an orthogonal array design to estimate all the main effects and interaction effects A X B and C X D.}{15}{4}{CO3}
\q{3}{Provide a concise summary of major Shainin Techniques.}{10}{2}{CO2}
\q{4}{Enumerate the four types of Signal-to-Ratios for the static problems}{10}{2}{CO2}
\end{cietable}
\end{iempaper}

% ---------- AS266TEA Fundamentals of Aerospace Engineering (CIE-3)
\begin{iempaperB}
\paperinfo{hdr=2,usn=0,rule=0,ayear={2025-2026 (Even Semester)},dept={AEROSPACE ENGINEERING},date={19\textsuperscript{th} June 2026},code={AS266TEA},sem={VI Semester},exam={CIE-3},maxmarks={60},duration={120 Min},
 title={FUNDAMENTALS OF AEROSPACE ENGINEERING},shownote=0}
\iemheadB
{\small Note:\ \ 1. Answer All Questions from both Part-A \& Part-B}\par
\hspace*{0.9cm}{\small 2. Part-A to be answered in the 1\textsuperscript{st} two pages of the answer booklet}\par\vspace{4pt}
\ptitlem{PART -- A}{(10 Marks)}
\begin{longtable}{|C{1.0cm}|p{11.7cm}|C{1.2cm}|C{1.2cm}|C{1.2cm}|}
\hline & \multicolumn{1}{c|}{\textbf{QUESTIONS}} & \textbf{M} & \textbf{CO} & \textbf{BL}\\\hline
\q{1.}{Which orbital element determines the shape of an orbit?}{01}{CO2}{L2}
\q{2.}{What is the type of orbit transfer, which is considered the most fuel-efficient orbital transfer method.}{01}{CO2}{L2}
\q{3.}{Bulkheads are generally oriented in the \blank[2cm] direction.}{01}{CO2}{L1}
\q{4.}{Truss structures are most efficient in resisting \blank[2cm] loads}{01}{CO2}{L2}
\q{5.}{Which structural member divides the fuselage into compartments?}{01}{CO2}{L1}
\q{6.}{Name any two structural members used in semi-monocoque construction.}{01}{CO2}{L1}
\q{7.}{A satellite has a perigee radius of 7000 km and an apogee radius of 9000 km. Find the semi-major axis.}{01}{CO2}{L3}
\q{8.}{An aircraft's airspeed indicator receives a pitot pressure of 145 kPa and a static pressure of 120 kPa. Determine the dynamic pressure acting on the aircraft.}{01}{CO4}{L3}
\q{9.}{A satellite moves from perigee to apogee. Does its velocity increase or decrease? Comment on the Gravitational potential Energy and Kinetic energy at Perigee and apogee.}{01}{CO4}{L4}
\q{10.}{An aircraft climbs from 2500 ft to 8500 ft in 6 minutes. Determine the average rate of climb.}{01}{CO4}{L3}
\end{longtable}
\ptitle{PART-B}
\begin{longtable}{|C{1.0cm}|p{11.7cm}|C{1.2cm}|C{1.2cm}|C{1.2cm}|}
\hline \textbf{Sl.\newline No.} & \multicolumn{1}{c|}{\textbf{Questions}} & \textbf{M} & \textbf{CO} & \textbf{BT}\\\hline
\q{1}{State and explain the 3 Kepler's Laws of planetary motion. Mathematically provide the expressions for the 2\textsuperscript{nd} and 3\textsuperscript{rd} Law.}{10}{CO2}{L2}
\q{2}{An artificial satellite is launched into a circular orbit at an altitude of 200 km above the Earth's surface.\newline
Assume:
\begin{bul}\item Radius of Earth, $R_E$ = 6378km\item Mass of Earth, $M_E$ = 5.97 $\times$ 10$^{24}$kg\item Universal gravitational constant, $G$ = 6.67 $\times$ 10$^{-11}$N$\cdot$m$^2$/kg$^2$\end{bul}
(a) Calculate the orbital velocity of the satellite.\newline
(b) Determine the orbital period of revolution of the satellite.\newline
(c) Derive the expression for the radius of a geosynchronous orbit and calculate its value assuming the Earth's rotational period is 86,164 s.\newline
(d) Calculate the escape velocity from the Earth's surface and compare it with the orbital velocity obtained in part (a).}{10}{CO3}{L3}
\q{3}{A spacecraft is initially in a circular Earth orbit at an altitude of 150 miles above the Earth's surface. It is required to transfer to another circular orbit at an altitude of 500 miles using a Hohmann transfer orbit.\newline
(a) Using the conservation of angular momentum, prove that
\par\vspace{2pt}\centerline{$R_pV_p=R_aV_a$}\par\vspace{2pt}
where $R_p$ and $R_a$ are the perigee and apogee radii of the transfer orbit, and $V_p$ and $V_a$ are the corresponding velocities.\newline
(b) Using conservation of energy, derive expressions for the velocity at perigee and apogee of an elliptical orbit:
\par\vspace{2pt}\centerline{$V_p=\sqrt{\dfrac{2GMR_a}{R_p(R_a+R_p)}}$\qquad $V_a=\sqrt{\dfrac{2GMR_p}{R_a(R_a+R_p)}}$}\par\vspace{2pt}
(c) Determine the velocity increments ($\Delta V_1$ and $\Delta V_2$) required to transfer the spacecraft from the 150-mile orbit to the 500-mile orbit.}{10}{CO3}{L3}
\q{4}{With a neat diagram of the Pitot-Static System, explain the working principles of the Airspeed Indicator (ASI), Altimeter, and Vertical Speed Indicator (VSI). Discuss the effect of pitot tube blockage and static port blockage on instrument indications}{10}{CO4}{L4}
\q{5}{With neat sketches, explain the construction and load-carrying mechanisms of Truss, Monocoque, and Semi-Monocoque aircraft structures. Compare them with respect to strength, weight, damage tolerance, and suitability for modern aircraft applications.}{10}{CO2}{L3}
\end{longtable}
\btfoot{|l|l|l|*{10}{C{0.85cm}|}}{\hline
\multirow{2}{*}{\shortstack{Marks\\Distribution}} & Particulars & & CO1 & CO2 & CO3 & CO4 & L1 & L2 & L3 & L4 & L5 & L6\\\cline{2-13}
 & Test & \shortstack{Max\\Marks} & 0 & 27 & 20 & 13 & 3 & 13 & 33 & 11 & -- & --\\\hline}
\end{iempaperB}

% ---------- BT266TEB Healthcare Analytics (Test - 3)
\begin{iempaperB}
\paperinfo{ayear={2025-2026 (Even Sem)},usn=0,rule=0,dept={BIOTECHNOLOGY},date={June, 2026},code={BT266TEB},sem={VI Semester},exam={Closed Book Offline Quiz and Test -- 3},maxmarks={10 + 50},duration={120 min},
 title={Healthcare Analytics (Institutional Elective)},shownote=0}
\vspace{-0.2cm}\centerline{Common to all branches}\par
\iemheadB
{\small Answer all questions in Part A \& Part B, Answer Part A in first two pages in sequence only}\par\vspace{2pt}
\ptitle{Part A}
\begin{xtable}{Questions}
\q{1}{What are sequencing adapters in Next-Generation Sequencing (NGS)?}{2}{2}{1}
\q{2}{What does a Phred quality score of 30 signify in DNA sequencing data?}{2}{2}{1}
\q{3}{Name any two commonly employed DNA enrichment methods used in NGS workflows.}{2}{2}{2}
\q{4}{What is meant by the term ``base calling'' in the context of Next-Generation Sequencing?}{2}{3}{2}
\q{5}{What is one significant limitation associated with read clipping during NGS data preprocessing?}{2}{2}{3}
\end{xtable}
\ptitle{Part B}
\begin{xtable}{Questions (Test)}
\q{1.}{Critically examine the role of quality assessment in NGS data analysis. How are Phred quality scores calculated, and what is their importance in ensuring reliable downstream bioinformatics analyses?}{10}{2}{3}
\q{2}{Discuss on the evolution of DNA sequencing technologies with special reference to the development of Sanger sequencing.}{10}{3}{3}
\q{3}{Elucidate on the sources and consequences of adapter and primer contamination in NGS datasets.}{10}{4}{4}
\q{4}{Compare first-generation and next-generation sequencing technologies in terms of methodology, scalability, accuracy, throughput, and applications in genomic research.}{10}{4}{4}
\q{5}{Discuss the implications for genomics research and personalized medicine.}{10}{2}{3}
\end{xtable}
\btfoot{|l|l|l|*{10}{C{0.9cm}|}}{\hline
\multirow{2}{*}{Marks Distribution} & Particulars & & CO1 & CO2 & CO3 & CO4 & L1 & L2 & L3 & L4 & L5 & L6\\\cline{2-13}
 & Test & Max Marks & & & 60 & 40 & & 40 & 20 & 40 & & \\\hline}
\par\vspace{4pt}\noindent{\itshape\small 1.\ \ \underline{(Include Nobel Prize, chain termination principle, Human Genome Project impact, etc.)}}
\end{iempaperB}

% ---------- CH266TEC Industrial Safety Engineering (Improvement Test)
\begin{xpaper}
\paperinfo{deptline={Department of Chemical Engineering},date={19.06.2026},exam={Improvement Test},maxmarks={60},sem={6},prog={UG},duration={2 Hrs},
 title={Industrial Safety Engineering},code={CH266TEC}}
\xhead
\begin{longtable}{|C{1.05cm}|p{12.2cm}|C{0.95cm}|C{0.95cm}|C{0.95cm}|}
\hline \textbf{S No} & \multicolumn{1}{c|}{\textbf{Questions}} & \textbf{M} & \textbf{BT} & \textbf{CO}\\\hline
\qpart{PART A}
\q{1}{Mention two employee responsibilities in process industry.}{2}{1}{2}
\q{2}{Write the types of body protection PPE kits and mention their use.}{2}{1}{4}
\q{3}{How many MIC storage tanks were installed at the Union Carbide Bhopal plant, and how many were designated for normal operation and emergency use?}{2}{1}{4}
\q{4}{What were direct casualties reported in Chernobyl disaster?}{2}{2}{4}
\q{5}{What is the third-year present value factor, up to three decimal places, if a 7\% risk-free rate of interest is used to calculate the Net Present Value?}{2}{3}{3}
\qpart{PART B}
\q{1}{Describe the potential incidents of eye and face related hazards. Classify and explain the eye and face PPE kits with sketches.}{10}{3}{4}
\q{2}{Mention the need of protective equipment for industrial works. Explain six-foot protection personal protective equipment's with sketches.}{10}{2}{4}
\q{3}{Explain briefly the process chemistry of union carbide Bhopal plant along with the process flow sheet.}{10}{3}{4}
\q{4}{The table below presents the possible net cash flows and their associated probabilities for Projects X and Y. Assuming a discount rate of 10\% and an initial investment of Rs. 5,000 for both projects, calculate the expected Net Present Value (NPV) for each project. Additionally, compute the variance, standard deviation, and coefficient of variation. Based on your analysis, indicate which project is preferable.
\qtabc{|c|c|c|c|c|}{\hline \multirow{2}{*}{Possible events} & \multicolumn{2}{c|}{Project X} & \multicolumn{2}{c|}{Project Y}\\\cline{2-5} & CF & P & CF & P\\\hline A & 4,000 & 0.1 & 12,000 & 0.10\\\hline B & 5,000 & 0.2 & 10,000 & 0.15\\\hline C & 6,000 & 0.4 & 8,000 & 0.50\\\hline D & 7,000 & 0.2 & 6,000 & 0.15\\\hline E & 8,000 & 0.1 & 4,000 & 0.10\\\hline}}{20}{4}{3}
\end{longtable}
\end{xpaper}

% ---------- CV266TEE Intelligent Transport Systems (Improvement CIE)
\begin{xpaper}
\paperinfo{deptline={Department of Civil Engineering},date={16-06-2026},exam={Improvement CIE},maxmarks={60},sem={VI},prog={UG},duration={2 Hrs},
 title={Intelligent Transport Systems},code={CV266TEE}}
\xhead
\ptitle{Quiz}
\begin{longtable}{|C{1.05cm}|p{12.2cm}|C{0.95cm}|C{0.95cm}|C{0.95cm}|}
\hline \textbf{S No} & \multicolumn{1}{c|}{\textbf{Questions}} & \textbf{M} & \textbf{BT} & \textbf{CO}\\\hline
\q{1.}{Explain the purpose of ITS evaluation in transportation projects.}{2}{2}{3}
\q{2.}{List the benefits provided by different ITS components in traffic management.}{2}{2}{4}
\q{3.}{Explain how ITS supports law enforcement in enforcing traffic rules and regulations.}{2}{2}{3}
\q{4.}{List the various funding options available for ITS implementation.}{2}{2}{4}
\q{5.}{List the importance of ITS standards testing before deployment.}{2}{2}{3}
\end{longtable}
\ptitle{Test}
\begin{longtable}{|C{1.05cm}|p{12.2cm}|C{0.95cm}|C{0.95cm}|C{0.95cm}|}
\hline \textbf{S No} & \multicolumn{1}{c|}{\textbf{Questions}} & \textbf{M} & \textbf{BT} & \textbf{CO}\\\hline
\q{1.}{Explain the various stages involved in ITS evaluation, including project selection, deployment tracking, impact assessment, and evaluation guidelines.}{10}{3}{3}
\q{2.}{Discuss the benefits of ITS components and explain how they contribute to improving transportation system performance.}{10}{3}{3}
\q{3.}{Explain the role of ITS in enhancing and supporting the enforcement of traffic rules and regulations with suitable examples.}{10}{3}{4}
\q{4.}{Explain the application areas of ITS standards and discuss the role of National Transportation Communications for ITS Protocol (NTCIP) in ITS systems.}{10}{3}{4}
\q{5.}{Discuss the role of ITS in the development of smart cities and explain how ITS applications improve urban mobility, safety, and sustainability.}{10}{3}{4}
\end{longtable}
\end{xpaper}

% ---------- CV266TEF Integrated Health Monitoring of Structures (Improvement CIE)
\begin{xpaper}
\paperinfo{deptline={Department of Civil Engineering},date={19.06.2026},exam={Improvement CIE},maxmarks={Q10 + T50},sem={VI},prog={UG},duration={2 Hrs},
 title={Integrated Health Monitoring of Structures (Institutional Elective)},code={CV266TEF}}
\xhead
{\small Note:}\par\vspace{-1pt}
\hspace*{0.5cm}{\footnotesize i) Answer all the questions.}\par
\hspace*{0.5cm}{\footnotesize ii) Use illustrative sketches to substantiate answers when necessary.}\par\vspace{3pt}
\ptitle{PART A}
\begin{longtable}{|C{1.05cm}|p{12.2cm}|C{0.95cm}|C{0.95cm}|C{0.95cm}|}
\hline \textbf{Sl. No.} & \multicolumn{1}{c|}{\textbf{Questions}} & \textbf{M} & \textbf{BT} & \textbf{CO}\\\hline
\q{1.}{List any two sensors used in Remote Structural Health Monitoring.}{02}{L1}{CO4}
\q{2.}{List any two limitations of Remote Structural Health Monitoring.}{02}{L1}{CO4}
\q{3.}{Discuss the role of IoT in Remote Structural Health Monitoring?}{02}{L2}{CO4}
\q{4.}{List the advantages of dynamic response methods.}{02}{L1}{CO3}
\q{5.}{Discuss the role of data acquisition systems in dam monitoring?}{02}{L2}{CO3}
\end{longtable}
\ptitle{PART B}
\begin{longtable}{|C{1.05cm}|p{12.2cm}|C{0.95cm}|C{0.95cm}|C{0.95cm}|}
\hline \textbf{Sl. No.} & \multicolumn{1}{c|}{\textbf{Questions}} & \textbf{M} & \textbf{BT} & \textbf{CO}\\\hline
\q{1.}{List and discuss the types of Dynamic field tests.}{10}{L2}{CO3}
\q{2.}{Discuss the advantages of Remote structural health monitoring.}{10}{L2}{CO4}
\q{3.}{Discuss the applications of Dynamic Field Testing.}{10}{L2}{CO3}
\q{4.}{Describe the various methods of Non-Destructive Evaluation used in structural health monitoring and condition assessment.}{10}{L2}{CO4}
\q{5.}{Explain the various dynamic response methods used for structural condition assessment.}{10}{L2}{CO3}
\end{longtable}
\par\vspace{4pt}\noindent{\bfseries Course outcomes:}\par
\btfoot{|l|l|l|*{10}{C{0.9cm}|}}{\hline
\multirow{2}{*}{Marks Distribution} & Particulars & & CO1 & CO2 & CO3 & CO4 & L1 & L2 & L3 & L4 & L5 & L6\\\cline{2-13}
 & Test & Max Marks & -- & -- & 34 & 26 & 06 & 54 & -- & -- & -- & --\\\hline}
\stars{***********}
\end{xpaper}

% ---------- CM266TEG Advanced Energy Storage for E-Mobility (Improvement CIE)
\begin{xpaper}
\paperinfo{deptline={Department of Chemistry},date={19-06-2026},exam={Improvement CIE},maxmarks={60 (50+10)},sem={6},prog={UG},duration={2Hrs (30+90 min)},
 title={ADVANCED ENERGY STORAGE FOR E-MOBILITY},code={CM266TEG}}
\xhead
{\small Instructions to candidate}\par\vspace{-1pt}
\hspace*{0.5cm}{\footnotesize i) \ All questions are compulsory\ \ ii) Part A questions should be answered in first two pages only}\par\vspace{3pt}
\ptitle{PART A}
\begin{longtable}{|C{1.05cm}|p{12.2cm}|C{0.95cm}|C{0.95cm}|C{0.95cm}|}
\hline \textbf{Sl NO} & \multicolumn{1}{c|}{\textbf{Questions}} & \textbf{M} & \textbf{BT} & \textbf{CO}\\\hline
\q{1}{State any one limitation of lithium-ion batteries in electric vehicles.}{1}{2}{1}
\q{2}{What is the electrolyte used in a lead-acid battery?}{1}{2}{1}
\q{3}{Zebra battery generally operated at high temperature. Justify.}{1}{3}{3}
\q{4}{Name any one anode materials used in sodium ion batteries.}{1}{2}{1}
\q{5}{Write the cathode reaction in sodium air batteries when aqueous electrolytes are used.}{1}{2}{4}
\q{6}{Give an example for asymmetric capacitor.}{1}{2}{1}
\q{7}{Justify the role of membrane separator in redox flow batteries.}{1}{3}{4}
\q{8}{What is regenerative braking?}{1}{2}{2}
\q{9}{Write any one limitation of solar cell hybrid vehicles.}{1}{3}{2}
\q{10}{Identify any one advantages of fuel cell hybrid vehicles.}{1}{2}{2}
\end{longtable}
\ptitle{PART B}
\begin{longtable}{|C{1.05cm}|p{12.2cm}|C{0.95cm}|C{0.95cm}|C{0.95cm}|}
\hline \textbf{S No} & \multicolumn{1}{c|}{\textbf{Questions}} & \textbf{M} & \textbf{BT} & \textbf{CO}\\\hline
\q{1}{With neat labelled diagram outline the construction and working of Nickel metal hydride battery with necessary reactions. List any three advantages and limitations.}{10}{2}{3}
\q{2}{Illustrate the construction and working of following batteries with respect to e mobility. i) Sodium--Sulphur battery ii) Types of sodium ion batteries with neat labelled diagram.}{10}{2}{1}
\q{3}{Explain the construction and working of Zebra battery along with its advantages and limitations. Comment on its efficiency over other batteries.}{10}{2}{4}
\q{4}{i) Outline the construction and working of pseudo type super capacitor. Why pseudo type supercapacitors are superior than that of EDLC?\newline
ii) Detail the working mechanism of battery supercapacitor hybrid electric vehicle along with flow chart.}{10}{4}{2}
\q{5}{Alkali electrolytes are not preferred in methanol oxygen fuel cell. Justify. Outline the construction and working of amorphous silicon solar cell and methanol oxygen fuel cell with neat labelled diagram.}{10}{2}{2}
\end{longtable}
\end{xpaper}

% ---------- ME266TES Automotive Mechatronics (Improvement CIE)
\begin{xpaper}
\paperinfo{deptline={DEPARTMENT OF MECHANICAL ENGINEERING},date={17-06-2026},exam={Improvement CIE},maxmarks={60},sem={VI},prog={UG},duration={2 Hrs},
 title={Automotive Mechatronics (Institutional Elective)},code={ME266TES}}
\xhead
\ptitle{QUIZ}
\begin{longtable}{|C{1.05cm}|p{12.2cm}|C{0.95cm}|C{0.95cm}|C{0.95cm}|}
\hline \textbf{\#} & \multicolumn{1}{c|}{\textbf{Question}} & \textbf{M} & \textbf{BT} & \textbf{CO}\\\hline
\q{1}{Telematics is a combination of telecommunications and \blank[2cm].}{1}{1}{1}
\q{2}{Radio communication in vehicles uses \blank[2cm] waves.}{1}{1}{4}
\q{3}{The speed of radio waves in free space is approximately \blank[1.5cm] $\times$ 10$^8$ m/s.}{1}{1}{4}
\q{4}{The sensor used to measure oxygen content in exhaust gases is called the \blank[2cm] sensor.}{1}{1}{1}
\q{5}{The crankshaft position sensor provides information about engine \blank[2cm].}{1}{1}{1}
\q{6}{The boost pressure sensor is mainly used in \blank[2cm] engines.}{1}{1}{1}
\q{7}{The coolant temperature sensor monitors the temperature of engine \blank[2cm].}{1}{1}{1}
\q{8}{MAF stands for Mass Air \blank[2cm].}{1}{1}{1}
\q{9}{The throttle position sensor measures the angle of the \blank[2cm] valve.}{1}{1}{1}
\q{10}{Rain sensors are commonly used for automatic \blank[2cm] operation.}{1}{1}{1}
\end{longtable}
\ptitle{TEST}
\begin{longtable}{|C{1.05cm}|p{12.2cm}|C{0.95cm}|C{0.95cm}|C{0.95cm}|}
\hline \textbf{\#} & \multicolumn{1}{c|}{\textbf{Question}} & \textbf{M} & \textbf{BT} & \textbf{CO}\\\hline
\q{1}{A connected vehicle transmits real-time location, speed, and engine health information to a remote service center. Explain the concept of telematics, radio transmission, signal path, and information exchange involved in this system.}{10}{3}{4}
\q{2}{Discuss the construction, working principle, applications, and advantages of oxygen sensors and crankshaft/camshaft sensors used in modern automobiles.}{10}{2}{1}
\q{3}{A turbocharged vehicle is experiencing poor acceleration and increased fuel consumption. Analyze the role of the boost pressure sensor and Hot Film Air Mass Flow Sensor in diagnosing the problem.}{10}{4}{1}
\q{4}{Explain the working principle, construction, and automotive applications of coolant temperature sensors, throttle position sensors, and rain/light sensors.}{10}{2}{1}
\q{5}{A modern vehicle exhibits erratic engine performance and intermittent communication faults between control modules. Discuss how telematics systems and automotive sensors assist in monitoring, diagnosis, and communication within the vehicle.}{10}{5}{4}
\end{longtable}
\end{xpaper}

% ---------- HS266TEW Applied Psychology for Engineers (CIE-III Improvement Test)
\begin{iempaper}
\paperinfo{shownote=0,date={19.06.2026},exam={CIE -- III -- Improvement Test},maxmarks={10 + 50},sem={VI},prog={UG- Institutional Elective},duration={2 hrs},
 title={Applied Psychology for Engineers},code={HS266TEW}}
\begin{center}{\fontsize{12}{14}\selectfont DEPARTMENT OF}\\[3pt]{\bfseries\fontsize{15.5}{18}\selectfont HUMANITIES AND SOCIAL SCIENCES}\end{center}\vspace{-3pt}
\noindent\begingroup\renewcommand{\arraystretch}{2.0}\bfseries\fontsize{11.5}{13}\selectfont
\begin{tabular}{|p{6.2cm}|C{5.0cm}|p{5.5cm}|}\hline
 Date \hspace{0.5cm}: \hspace{0.2cm}19.06.2026 & CIE -- III --Improvement Test & Max. Marks \hspace{0.3cm}: \hspace{0.2cm}10 + 50\\\hline
 Semester \hspace{0.1cm}: \hspace{0.2cm}VI & UG- Institutional Elective & Duration \hspace{0.55cm}: \hspace{0.2cm}2 hrs\\\hline
 \multicolumn{2}{|p{11.2cm}|}{Course Title: Applied Psychology for Engineers} & Course Code \hspace{0.1cm}: \hspace{0.2cm}HS266TEW\\\hline
\end{tabular}\endgroup\par\vspace{8pt}
\begin{cietable}{M}{BT}{CO}
\qpart{Part A \hfill 10 Marks}
\q{1.1}{Define the trait approach to personality.}{2}{L4}{CO4}
\q{1.2}{Mention the psychological impact of stress.}{2}{L2}{CO1}
\q{1.3}{Mention the importance of learning in human behaviour.}{2}{L3}{CO4}
\q{1.4}{What is counselling? Mention the need for counselling.}{2}{L3}{CO3}
\q{1.5}{What is the role of an organizational psychologist in improving employee performance?}{2}{L2}{CO2}
\qpart{Part -- B \hfill 50 Marks}
\q{2}{A young professional is highly motivated to develop personal skills, achieve goals, and reach his full potential. He believes in self-growth, making independent choices, and finding meaning in life despite challenges. Using the humanistic approach to personality, explain how personality development occurs in this individual. Discuss the key concepts of the humanistic approach.}{10}{L1}{CO1}
\q{3}{A psychologist is assessing a client who finds it difficult to express personal feelings, emotions, and inner conflicts directly. The psychologist uses ambiguous images and incomplete scenarios to understand the client's underlying thoughts and personality traits. Based on the given case, explain how projective techniques help the psychologist in understanding the client's personality and discuss the different types of projective techniques used in personality assessment.}{10}{L2}{CO2}
\q{4}{A student feels anxious whenever entering the examination hall because of previous experiences of stress and fear during exams. Over time, the examination hall itself becomes associated with anxiety, even before the exam begins. Using the concept of classical conditioning, explain how the student developed anxiety towards the examination hall. Identify the components of classical conditioning involved in this situation.}{10}{L3}{CO2}
\q{5}{An organization observes that some employees are experiencing stress, low motivation, and difficulty in adapting to workplace challenges. To support employees and improve their performance, the organization introduces counselling services. Discuss the different approaches of counselling used to support employees.}{10}{L4}{CO3}
\q{6}{The company is recruiting new employees and wants to select candidates with suitable skills and personality traits. After the selection process, the company plans to provide training to improve employee performance and workplace skills. Explain the steps involved in the employee selection and training process.}{10}{L2}{CO4}
\end{cietable}
\btfoot{|l|l|l|*{11}{C{0.8cm}|}}{\hline
\multirow{2}{*}{\shortstack{Marks\\Distribution}} & Particulars & & CO1 & CO2 & CO3 & CO4 & CO5 & L1 & L2 & L3 & L4 & L5 & L6\\\cline{2-14}
 & Quiz & Max Marks & 2 & 4 & 2 & 2 & -- & 2 & 2 & 2 & 4 & -- & --\\\cline{2-14}
 & Test & Max Marks & 10 & 20 & 10 & 10 & -- & 10 & 20 & 20 & -- & -- & --\\\hline}
\end{iempaper}

% ---------- IM266TEQ Elements of Financial Management (Improvement CIE)
\begin{iempaper}
\paperinfo{shownote=0,date={19\textsuperscript{th} June 2026},exam={Improvement CIE},maxmarks={10 + 50},sem={VI Semester},duration={120 Mins},
 title={ELEMENTS OF FINANCIAL MANAGEMNT},code={IM266TEQ}}
\iemheadB
{\small Note:}\par\vspace{-1pt}
\hspace*{0.5cm}{\footnotesize 1. \hspace{0.1cm}Answer all the Questions.}\par
\hspace*{0.5cm}{\footnotesize 2. \hspace{0.1cm}Avoid verbose answers and base them on your general knowledge. You should provide evidence of you having acquired the Subject Knowledge related to the course.}\par\vspace{3pt}
\ptitle{PART -- A}
\begin{longtable}{|C{1.05cm}|p{12.2cm}|C{0.95cm}|C{0.95cm}|C{0.95cm}|}
\hline \textbf{Sl.\newline No.} & \multicolumn{1}{c|}{\textbf{Questions}} & \textbf{M} & \textbf{BT} & \textbf{CO}\\\hline
\q{1}{State the decision rule of NPV}{01}{L$_2$}{CO3}
\q{2}{Which technique of capital budgeting ignores cash flows after payback?}{01}{L$_2$}{CO3}
\q{3}{Which source of long-term financing increases financial risk?}{01}{L$_2$}{CO4}
\q{4}{Which regulator in the Indian Financial System approves IPOs?}{01}{L$_1$}{CO1}
\q{5.}{Which market of the two broad categories related to buying and selling of securities deals with resale of securities?}{01}{L$_1$}{CO1}
\q{5.}{What is an IPO?}{01}{L$_1$}{CO1}
\q{6.}{What are Retained Earnings?}{01}{L$_1$}{CO3}
\q{7.}{What is Permanent Working Capital?}{01}{L$_2$}{CO3}
\q{8.}{Define Conservative Policy in the context of working capital policy?}{01}{L$_2$}{CO3}
\q{9.}{Which technique of capital budgeting may give multiple rates?}{01}{L$_2$}{CO3}
\q{10.}{Which source of long-term financing has a tax shield?}{01}{L$_2$}{CO4}
\end{longtable}
\ptitle{PART -- B}
\begin{longtable}{|C{1.05cm}|p{12.2cm}|C{0.95cm}|C{0.95cm}|C{0.95cm}|}
\hline \textbf{Sl.\newline No.} & \multicolumn{1}{c|}{\textbf{Questions}} & \textbf{M} & \textbf{BT} & \textbf{CO}\\\hline
\q{1.a}{Based on the Information Provided below\newline
Project A: Cost \rupee 1,00,000; inflows \rupee 40,000 per year for 4 years.\newline
Project B: Cost \rupee 1,00,000; inflows \rupee 20,000 for 2 years, then \rupee 70,000 for 2 years.\newline
Discount rate = 10\%. Which of the projects should be selected based on the NPV decision rule? Rework the problem based on the IRR Rule? Discuss NPV vs IRR Conflict.}{10}{L$_4$}{CO3}
\q{2a.}{Mention the major sources of long-term finance with their advantages and limitations.}{06}{L$_2$}{CO4}
\q{b.}{Explain the advantages and limitations of debentures as a source of finance.}{04}{L$_2$}{CO4}
\q{3a.}{Discuss the core ideas and concepts related to financing mix for working capital.}{06}{L$_3$}{CO3}
\q{b.}{What is book building in IPOs?}{04}{L$_2$}{CO1}
\q{4a.}{Discuss the determinants of working capital in manufacturing Vis-\`a-vis Service firms.}{05}{L$_2$}{CO3}
\q{b.}{Given Cost = \rupee 3,00,000. Cash inflows: \rupee 1,00,000 per year for 4 years.\newline
Discount rates: 10\% and 15\%. PV factors (10\%, 4 years) = 3.17; (15\%, 4 years) = 2.855. Compute the NPV using Discount Rates of 10\% and 15\%.}{05}{L$_3$}{CO3}
\q{5a.}{Explain the trade-off between Liquidity vs. profitability in financial management and working capital policy.}{05}{L$_3$}{CO3}
\q{b.}{Explain the role of the secondary market in price discovery.}{05}{L$_3$}{CO1}
\end{longtable}
\end{iempaper}

% ---------- HS266TEY Universal Human Values - III (Improvement Test)
\begin{xpaper}
\paperinfo{deptline={},date={19 June 2026},exam={IMPROVEMENT TEST},maxmarks={50+10},sem={VI},prog={},duration={90+20 Min},
 title={Universal Human Values -- III},code={HS266TEY}}
\par\vspace{2pt}\centerline{\small Department of Civil Engineering}\par\vspace{6pt}
\noindent\begingroup\renewcommand{\arraystretch}{1.5}\bfseries\fontsize{11}{13}\selectfont
\begin{tabular}{|p{3.2cm}|p{4.8cm}|C{3.6cm}|C{3.6cm}|}\hline
 Date & 19 June 2026 & Maximum Marks & 50+10\\\hline
 Course Code & HS266TEY & Duration & 90+20 Min\\\hline
 Sem & VI & \multicolumn{2}{l|}{IMPROVEMENT TEST}\\\hline
 \multicolumn{4}{|c|}{Universal Human Values -- III}\\\hline
\end{tabular}\endgroup\par\vspace{8pt}
\begin{longtable}{|C{1.05cm}|p{12.2cm}|C{0.95cm}|C{0.95cm}|C{0.95cm}|}
\hline \textbf{Sl. No.} & \multicolumn{1}{c|}{\textbf{Quiz}} & \textbf{M} & \textbf{BT} & \textbf{CO}\\\hline
\q{1}{Nature is made up of \blank[1.6cm] and \blank[1.6cm].}{02}{1}{2}
\q{2}{Co-existence is ever-present, \blank[1.6cm] \& \blank[1.6cm].}{02}{1}{2}
\q{3}{Which are the four orders in which coexistence is expressed?}{02}{1}{1}
\q{4}{Unit is \blank[1.6cm] being in space. So, it has definite \blank[1.6cm] in space.}{02}{2}{2}
\q{5}{What is the human goal at the level of self?}{02}{2}{1}
\end{longtable}
\par\vspace{6pt}
\begin{longtable}{|C{1.05cm}|p{12.2cm}|C{0.95cm}|C{0.95cm}|C{0.95cm}|}
\hline \textbf{Sl. No.} & \multicolumn{1}{c|}{\textbf{Test}} & \textbf{M} & \textbf{BT} & \textbf{CO}\\\hline
\q{1}{Explain the distinct roles of self-exploration, self-awareness, Self-evaluation in the process of inner evaluation.}{10}{2}{2}
\q{2}{Examine the statement `` when realization, understanding and contemplation in the self, human behavior and work naturally become harmonious''}{10}{2}{3}
\q{3}{How does achieving harmony inside the self naturally lead to the smooth functioning and nurturing of the physical body.}{10}{3}{2}
\q{4}{Differentiate between the self and body based on their needs and activities.}{10}{3}{3}
\q{5}{What is the role of human being in this ever-expressing co-existence?}{10}{2}{4}
\end{longtable}
\par\vspace{4pt}\noindent{\bfseries Course outcomes:}\par
\noindent\begin{tabular}{|C{1.3cm}|p{15cm}|}\hline CO 1 & Understand the basic human aspiration with program of its fulfillment and meaning of resolution in the complete expanse of human living\\\hline CO 2 & Understand human being in depth and see how self is central to human being\\\hline CO 3 & Understand existence in depth and see how coexistence is central to existence\\\hline CO 4 & Understand human conduct and the holistic way of living leading to human tradition\\\hline\end{tabular}
\btfoot{|l|l|l|*{10}{C{0.9cm}|}}{\hline \multirow{2}{*}{Marks Distribution} & Particulars & & CO1 & CO2 & CO3 & CO4 & L1 & L2 & L3 & L4 & L5 & L6\\\cline{2-13} & Test & Max Marks & 04 & 26 & 20 & 10 & 06 & 34 & 20 & -- & -- & --\\\hline}
\stars{***********}
\end{xpaper}
